\begin{proposition} \label{proposition:derivative_of_natural_logarithm_is_reciprocal}
The \CrefAndHyperrefIfExist{definition:real_natural_logarithm_function_as_inverse_of_real_exponential_function}{natural logarithm} $\ln : (0, \infty) \to \mathbb{R}$ is \CrefAndHyperrefIfExist{definition:derivative_of_a_real_valued_function_on_a_real_interval}{differentiable} at every point $x > 0$, and
$$\frac{d}{dx} \ln(x) = \frac{1}{x}$$
\end{proposition}
\begin{proof}
    Since $\ln(x)$ is the \CrefAndHyperrefIfExist{definition:injective_surjective_bijective_map_of_sets}{inverse} of \CrefAndHyperrefIfExist{definition:real_exponential_function_with_eulers_number_as_base}{$\exp$}, we have 
    $$x = \exp(\ln(x))$$
    for all $x \in (0,\infty)$. By the \CrefAndHyperrefIfExist{proposition:sum_rule_product_rule_quotient_rule_chain_rule_for_real_derivatives_of_real_valued_functions_on_real_intervals}{chain rule}, this yields
    $$1 = \frac{d}{dx} x = \frac{d}{dx} (\exp(\ln(x))) = \frac{d}{d y} \exp(y) \cdot \frac{d}{dx} \ln(x)$$
    where $y = \ln(x)$. Therefore, 
    $$1 = \exp(y) \cdot \frac{d}{dx} \ln(x),$$
    \CrefIfExists{proposition:derivative_of_exponential_function_is_itself}
    so 
    $$\frac{d}{dx} \ln(x) = \frac{1}{\exp(y)} = \frac{1}{\exp(ln(x))} = \frac{1}{x}$$
    as desired.
\end{proof}
